\NSPage{127}%
\begin{proof}
\begin{EESegment}{EE-TGT-P127-001}
A nonzero linear combination of \NSi{NS-F-P127-001}{m} distinct powers can have at
most \NSi{NS-F-P127-002}{m-1} distinct positive zeros. To see this, divide by the
smallest power. If the quotient had \NSi{NS-F-P127-003}{m} zeros, Rolle's theorem
would give \NSi{NS-F-P127-004}{m-1} zeros of its derivative. That would contradict the
induction hypothesis for the other \NSi{NS-F-P127-005}{m-1} powers. It follows that
\NSi{NS-F-P127-006}{\det[x_j^{\alpha_i}]} is nonzero whenever
\NSi{NS-F-P127-007}{0<x_1<\cdots<x_m}, and its sign stays the same throughout that
connected set. Since the determinant is multilinear,
\NSFormula{NS-F-P127-008}{display}{none}
\[
\det B
=\int_{I_1\times\cdots\times I_m}\det[x_j^{\alpha_i}]
  \prod_j\beta_j(x_j)\,dx_1\cdots dx_m\ne0.
\]
The integrand has one sign, and it is nonzero on a set of positive measure. The change of
variables \NSi{NS-F-P127-009}{x=e^y} proves the logarithmic version. Compactness,
together with differentiation of \NSi{NS-F-P127-010}{B^{-1}B=\Id}, proves the final
statement. That statement does not give a uniform bound for the inverse when two exponents
approach each other.
\end{EESegment}
\end{proof}

\begin{EESegment}{EE-TGT-P127-002}
Some moment weights approach each other as \NSi{NS-F-P127-011}{\lambda} tends to
zero. The corrections have to stay small enough to preserve the cone, so the proof needs to
know how much this costs in the inverse bound. Take identical translated bumps
\NSi{NS-F-P127-012}{\beta(y-c_1)} and
\NSi{NS-F-P127-013}{\beta(y-c_2)}, and set
\NSi{NS-F-P127-014}{\Delta=c_2-c_1>0} and
\NSi{NS-F-P127-015}{b(s)=\int e^{st}\beta(t)\,dt}. Dividing row
\NSi{NS-F-P127-016}{i} by \NSi{NS-F-P127-017}{e^{s_i c_1}} gives
\NSFormula{NS-F-P127-018}{display}{A.1}
\begin{equation}
B=
\begin{pmatrix}
b(s_1)&e^{s_1\Delta}b(s_1)\\
b(s_2)&e^{s_2\Delta}b(s_2)
\end{pmatrix},
\qquad
\det B=b(s_1)b(s_2)(e^{s_2\Delta}-e^{s_1\Delta}).
\tag{A.1}\label{eq:A.1}
\end{equation}
Fix the bump and fix \NSi{NS-F-P127-019}{\Delta>0}. Then
\NSi{NS-F-P127-020}{\lVert B^{-1}\rVert=O(\lambda^{-1})} when
\NSi{NS-F-P127-021}{s_1-s_2=\lambda} and the slopes stay in a fixed compact interval.
The inverse stays
bounded if the limiting separation of the slopes is nonzero. The same
\NSi{NS-F-P127-022}{O(\lambda^{-1})} bound holds for moment matrices with weights
\NSi{NS-F-P127-023}{1,x^{-\lambda}} and bumps on ordered disjoint intervals. This
follows by expanding
\NSi{NS-F-P127-024}{x^{-\lambda}=1-\lambda\log x+O(\lambda^2)}, because the two
bump averages of \NSi{NS-F-P127-025}{\log x} stay strictly separated.
\end{EESegment}

\begin{EESegment}{EE-TGT-P127-003}
The rank calculation handles the linearized moment map. The actual moments are quadratic
in the velocity profiles, so the correction coefficients have to solve a finite system with a
quadratic remainder. The next estimate gives a small solution and also controls its
parameter derivatives. It does not require all orders of derivatives to be small at the same
time.
\end{EESegment}

\begin{lemma}\label{lem:A.2}
\begin{EESegment}{EE-TGT-P127-004}
Let \NSi{NS-F-P127-026}{K=[-1,1]}. Suppose a finite family of correction coefficients
\NSi{NS-F-P127-027}{c\in\R^m} changes the required moments in the following way. First
take the specified linear combinations of the moments, and divide each resulting component
by its specified nonzero factor. The change is then
\NSFormula{NS-F-P127-028}{display}{A.2}
\begin{equation}
F_\eta(c)=B(\eta)c+\mathcal Q_\eta(c,c),
\qquad \eta\in K.
\tag{A.2}\label{eq:A.2}
\end{equation}
Here \NSi{NS-F-P127-029}{B} and the bilinear map
\NSi{NS-F-P127-030}{\mathcal Q} are smooth,
\NSi{NS-F-P127-031}{B} is invertible on \NSi{NS-F-P127-032}{K}, and the given
moment discrepancy is \NSi{NS-F-P127-033}{d\in C^\infty(K;\R^m)}. Write
\NSi{NS-F-P127-034}{\beta_0=\sup_K\lVert B^{-1}\rVert},
\NSi{NS-F-P127-035}{\kappa_0=\sup_K\lVert\mathcal Q\rVert}, and
\NSi{NS-F-P127-036}{d_0=\lVert d\rVert_{C^0(K)}}. If
\NSFormula{NS-F-P127-037}{display}{none}
\[
8\beta_0^2\kappa_0d_0\le1,
\]
there is exactly one solution of \NSi{NS-F-P127-038}{F_\eta(c)=d(\eta)} in the ball
\NSi{NS-F-P127-039}{\lVert c\rVert_{C^0}\le2\beta_0d_0}. This solution is smooth in
\NSi{NS-F-P127-040}{\eta}, and iterating
\NSi{NS-F-P127-041}{c\mapsto B(\eta)^{-1}(d(\eta)-\mathcal Q_\eta(c,c))} from zero
selects it. If \NSi{NS-F-P127-042}{\mathcal Q=0}, no smallness condition is needed.

Now choose any finite \NSi{NS-F-P127-043}{k} in advance. Let
\NSi{NS-F-P127-044}{\beta_k} be the operator norm for multiplication by
\NSi{NS-F-P127-045}{B^{-1}} on \NSi{NS-F-P127-046}{C^k(K)}, and let
\NSi{NS-F-P127-047}{\kappa_k} be the corresponding bilinear norm of
\NSi{NS-F-P127-048}{\mathcal Q}. If
\NSi{NS-F-P127-049}{8\beta_k^2\kappa_k\lVert d\rVert_{C^k}\le1} as well, the same
solution branch satisfies
\NSFormula{NS-F-P127-050}{display}{A.3}
\begin{equation}
\lVert c\rVert_{C^k}\le2\beta_k\lVert d\rVert_{C^k}.
\tag{A.3}\label{eq:A.3}
\end{equation}
\end{EESegment}
\end{lemma}

\begin{proof}
\begin{EESegment}{EE-TGT-P127-005}
Set \NSi{NS-F-P127-051}{r=2\beta_0d_0}. The map
\NSi{NS-F-P127-052}{c\mapsto B^{-1}(d-\mathcal Q(c,c))} takes the closed
\NSi{NS-F-P127-053}{C^0} ball of radius \NSi{NS-F-P127-054}{r} into itself. Its
Lipschitz constant on that ball is at most
\NSi{NS-F-P127-055}{2\beta_0\kappa_0r\le1/2}. The same inequality makes
\NSi{NS-F-P127-056}{B+D_c\mathcal Q(c,c)} invertible. The pointwise implicit function
 theorem then gives smoothness, including one-sided derivatives at the endpoints. Apply the
 same contraction argument to
\end{EESegment}
\NSPage{128}%
\begin{EESegment}{EE-TGT-P128-001}
\NSi{NS-F-P128-001}{C^k(K)}, the space of functions on the compact interval whose derivatives through the stated order
are continuous. This space is complete in its usual norm, and pointwise multiplication makes it an algebra. That is the
Banach algebra used here, and the argument gives \eqref{eq:A.3}. The iterations are the same ones as
in \NSi{NS-F-P128-002}{C^0}, so they have the same limit. Every other fixed derivative
is finite by implicit differentiation, with a constant that may depend on the corresponding
data. Any extra parameter in a fixed compact set, such as the pulse amplitude below, is
handled in exactly the same way.
\end{EESegment}
\end{proof}

\begin{EESegment}{EE-TGT-P128-002}
The size of these correction estimates depends on the radial scale of the support. For
example, if \NSi{NS-F-P128-003}{\delta V=v_*\sum_j c_j(\eta)\beta_j(X/\rho)}, then
\NSFormula{NS-F-P128-004}{display}{none}
\[
\lVert\partial_X^r\partial_\eta^k\delta V\rVert_\infty
\le C_rv_*\rho^{-r}\lVert c\rVert_{C^k}.
\]
For bumps whose width in \NSi{NS-F-P128-005}{\log X} stays fixed, the matching estimate
uses \NSi{NS-F-P128-006}{(X\partial_X)^r}. Small coefficients therefore preserve any
fixed finite collection of strict inequalities involving field derivatives on that patch. This
estimate applies to the moment-correction bumps. The rapid modulation that comes before a
correction has separate estimates for its radial derivatives.
\end{EESegment}

\begin{EESegment}{EE-TGT-P128-003}
For the five cumulative integrals in \eqref{eq:4.15}, the linearized system on a power-law
interval splits into two axial corrections and three azimuthal corrections. The next
corollary identifies their weights so that the rank calculation and the small-solution result
above can be applied.
\end{EESegment}

\begin{corollary}\label{cor:A.3}
\begin{EESegment}{EE-TGT-P128-004}
Use the five radial moment functions \NSi{NS-F-P128-007}{(M,I,J,S,C_p)} from
\eqref{eq:4.15}, where \NSi{NS-F-P128-008}{C_p} is the pressure increment. Under the
rescaling \NSi{NS-F-P128-009}{X=\rho x}, their normalizing factors are
\NSFormula{NS-F-P128-010}{display}{A.4}
\begin{equation}
(\rho,\rho^{3/2},\rho^{3/2},\rho,1),
\qquad H=\rho^{1/2}\sqrt{2x}E.
\tag{A.4}\label{eq:A.4}
\end{equation}
On a fixed correction interval, suppose
\NSi{NS-F-P128-011}{U_0=u_c(\eta)} and
\NSi{NS-F-P128-012}{E_0=e_*f_*(\eta)x^\alpha}, with
\NSi{NS-F-P128-013}{e_*>0} and smooth \NSi{NS-F-P128-014}{f_*>0}. Two additive
\NSi{NS-F-P128-015}{U} bumps and three additive \NSi{NS-F-P128-016}{E} bumps give an
invertible Jacobian for the normalized five-moment map as long as
\NSi{NS-F-P128-017}{\alpha\notin\{-1/2,1/2,3/2\}}. The moment changes have the exact
quadratic form \eqref{eq:A.2}. When \NSi{NS-F-P128-018}{U_0=0}, the three corrections
to \NSi{NS-F-P128-019}{E} can prescribe changes in
\NSi{NS-F-P128-020}{(I,S,C_p)} while leaving \NSi{NS-F-P128-021}{M,J} unchanged.
\end{EESegment}
\end{corollary}

\begin{proof}
\begin{EESegment}{EE-TGT-P128-005}
In the normalized rows, replace \NSi{NS-F-P128-022}{J} by
\NSi{NS-F-P128-023}{J-u_cI} and \NSi{NS-F-P128-024}{S} by
\NSi{NS-F-P128-025}{S-2u_cM}, keeping \NSi{NS-F-P128-026}{u_c} fixed at its base
value. The differential of \NSi{NS-F-P128-027}{J-u_cI} is
\NSi{NS-F-P128-028}{\int\sqrt{2x}E_0\,\delta U\,dx}, and the differential of
\NSi{NS-F-P128-029}{S-2u_cM} is
\NSi{NS-F-P128-030}{-\int E_0\,\delta E\,dx}. Together with the other three rows, these
give a matrix with two blocks. Apart from nonzero normalizing factors, their weights are
\NSFormula{NS-F-P128-031}{display}{none}
\[
\begin{array}{c|ccc}
U\text{ block}&1,&x^{\alpha+1/2}&\\
E\text{ block}&x^{1/2},&x^\alpha,&x^{\alpha-1}.
\end{array}
\]
The factors \NSi{NS-F-P128-032}{e_*,f_*} and
\NSi{NS-F-P128-033}{\sqrt2} remain in the normalization of the moment components. Their
inverses and every fixed parameter derivative stay bounded. The condition on
\NSi{NS-F-P128-034}{\alpha} says exactly that the powers in each block are distinct, so
Lemma~\ref{lem:A.1} applies. The original functionals have degree at most two in
\NSi{NS-F-P128-035}{U,E}, which makes every remaining term quadratic. In particular,
the row operations do not assume that the corrected \NSi{NS-F-P128-036}{U} stays
constant.
\end{EESegment}
\end{proof}

\begin{EESegment}{EE-TGT-P128-006}
For the axis moment correction \NSi{NS-F-P128-037}{\alpha=1/10}, the two blocks use
the powers \NSi{NS-F-P128-038}{(0,3/5)} and
\NSi{NS-F-P128-039}{(1/2,1/10,-9/10)}. On the intermediate power-law interval,
\NSi{NS-F-P128-040}{\alpha=-1/2-\lambda} gives
\NSi{NS-F-P128-041}{(0,-\lambda)} and
\NSi{NS-F-P128-042}{(1/2,-1/2-\lambda,-3/2-\lambda)}. The first of these last two
blocks can bring in an inverse bound of \NSi{NS-F-P128-043}{\lambda^{-1}}. The value of
\NSi{NS-F-P128-044}{\lambda} is fixed before the later radial frequency is chosen. These
normalized inverse bounds, together with the exact restoration of matching data in
Lemma~\ref{lem:4.4}, are what both corrections use.
\end{EESegment}

\NSPage{129}%
\begin{EESegment}{EE-TGT-P129-001}
\subsection{Constructing and matching the outer radial profile.}\label{sec:A.2}
This subsection gives a profile that connects an inner reference power law to a purely
azimuthal exterior power law \NSi{NS-F-P129-001}{E_{\mathrm{pow}}}. Its inner data will
be prescribed in \eqref{eq:A.7}. The intervals in between will set the required total
moments in \eqref{eq:A.8} while keeping the stress-cone inequalities. Four subintervals
will be kept free for the later profile, heat, background, and mean corrections; they are
listed after \eqref{eq:A.9}. First come the radial pieces. Their free amplitudes are chosen
and their stated properties checked afterward, in Sections~\ref{sec:A.3} and
\ref{sec:A.5}.
\end{EESegment}

\begin{EESegment}{EE-TGT-P129-002}
Fix, once and for all, the smooth step
\NSFormula{NS-F-P129-002}{display}{A.5}
\begin{equation}
\sigma(y)=\frac{e^{-1/y^2}}{e^{-1/y^2}+e^{-1/(1-y)^2}}
\quad(0<y<1),
\qquad \sigma=0\ (y\le0),
\qquad \sigma=1\ (y\ge1).
\tag{A.5}\label{eq:A.5}
\end{equation}
The positive moment-correction bumps are rescaled copies of
\NSi{NS-F-P129-003}{\sigma'} placed on ordered, disjoint subintervals of the chosen
patch. Each stage uses its own logarithmic radial coordinate
\NSi{NS-F-P129-004}{y}. Its origin is the start of that stage, and
\NSi{NS-F-P129-005}{d/dy=\mathcal D_X=X\partial_X}.
\end{EESegment}

\begin{EESegment}{EE-TGT-P129-003}
The parameters for the next profile are fixed in this order:
\NSFormula{NS-F-P129-006}{display}{A.6}
\begin{equation}
M_d,
\qquad T_d=e^{M_d}+10,
\qquad P_*>e^{T_d},
\qquad 0<\lambda\ll1,
\qquad 0<h\ll\min\{\lambda,e^{-T_d}\}
\tag{A.6}\label{eq:A.6}
\end{equation}
Here \NSi{NS-F-P129-007}{\ll} has the smallness meaning from
Definition~\ref{def:3.3}: choose \NSi{NS-F-P129-008}{\lambda} small enough after
fixing \NSi{NS-F-P129-009}{M_d,T_d,P_*}, and after that choose
\NSi{NS-F-P129-010}{h} small enough compared with both
\NSi{NS-F-P129-011}{\lambda} and \NSi{NS-F-P129-012}{e^{-T_d}}. The large radius
\NSi{NS-F-P129-013}{X_R} is chosen after all of them.
\end{EESegment}

\begin{EESegment}{EE-TGT-P129-004}
On the reference inner interval, write \NSi{NS-F-P129-014}{x=X/X_R} and
\NSi{NS-F-P129-015}{y=\log x}, and prescribe
\NSFormula{NS-F-P129-016}{display}{A.7}
\begin{equation}
U=4\eta,
\qquad E=P_*f(\eta)e^{y/10},
\qquad f=(1+\eta^2)^{-1}=e^{-J_0},
\qquad J_0=\log(1+\eta^2),
\qquad y\le0.
\tag{A.7}\label{eq:A.7}
\end{equation}
This ideal region is only a temporary choice for the outer radial profile; Section~\ref{sec:B}
replaces its inner part with a regular axis. The goal outside is the power
\NSi{NS-F-P129-017}{E_{\mathrm{pow}}=c_\infty X^{-A}}, where
\NSi{NS-F-P129-018}{c_\infty>0} does not depend on
\NSi{NS-F-P129-019}{\eta}, together with the moment identities
\NSFormula{NS-F-P129-020}{display}{A.8}
\begin{equation}
\int_0^\infty(U^2-E^2/2)\,dX=0,
\qquad
\int_0^\infty(H-H_{\mathrm{pow}})\,dX=0,
\qquad
H_{\mathrm{pow}}=\sqrt{2X}E_{\mathrm{pow}}.
\tag{A.8}\label{eq:A.8}
\end{equation}
These two identities remove the two integration constants that would otherwise remain in
the exterior stress. The profile is now built one stage at a time.
\end{EESegment}

\begin{EESegment}{EE-TGT-P129-005}
Recall that \NSi{NS-F-P129-021}{l=X\partial_X\log H}. Choosing
\NSi{NS-F-P129-022}{l} is therefore the same as choosing
\NSi{NS-F-P129-023}{(\log E)'=l-1/2} and integrating it.
\end{EESegment}

\begin{EESegment}{EE-TGT-P129-006}
\noindent\textbf{Reducing the axial component to zero.} Start at
\NSi{NS-F-P129-024}{x=1}. Keep \NSi{NS-F-P129-025}{U=4\eta}, and for one unit of
logarithmic radial length use \NSi{NS-F-P129-026}{l=(3/5)(1-\sigma(y))}. On the next
interval, set
\NSFormula{NS-F-P129-027}{display}{none}
\[
l=0,
\qquad U=k(y)\eta,
\qquad k(y)=4[1-\sigma(\log(1+y)/M_d)],
\qquad 0\le y\le T_d.
\]
Here \NSi{NS-F-P129-028}{|k'|\le e_d/(1+y)}, where
\NSi{NS-F-P129-029}{e_d=4\lVert\sigma'\rVert_\infty/M_d}. The function
\NSi{NS-F-P129-030}{k} is zero for the last eleven units of this stage. Let
\NSi{NS-F-P129-031}{P_1\asymp P_*} be the value of
\NSi{NS-F-P129-032}{E/f} at the start of the axial transition. Then
\NSi{NS-F-P129-033}{E=P_1fe^{-y/2}} throughout that interval.
\end{EESegment}

\begin{EESegment}{EE-TGT-P129-007}
\noindent\textbf{The intermediate power-law interval.} Set
\NSi{NS-F-P129-034}{U=0}. Use \NSi{NS-F-P129-035}{l=-\lambda\sigma(y)} on a unit
interval, then keep \NSi{NS-F-P129-036}{l=-\lambda} for
\NSFormula{NS-F-P129-037}{display}{A.9}
\begin{equation}
T_w=60\log(1/\lambda).
\tag{A.9}\label{eq:A.9}
\end{equation}
In the local logarithmic radial coordinate on this interval, reserve
\NSi{NS-F-P129-038}{(T_w-25,T_w-20)},
\NSi{NS-F-P129-039}{(T_w-20,T_w-15)},
\NSi{NS-F-P129-040}{(T_w-14,T_w-9)}, and
\NSi{NS-F-P129-041}{(T_w-8,T_w-3)}. In that order, they will hold the correction after
the cone is realized, the heat compensation, the moment conditions for higher-order
background corrections, and the moment conditions for the phase-averaged corrections.
Choose \NSi{NS-F-P129-042}{\lambda} small enough that all four patches lie inside this
interval.
\end{EESegment}

\NSPage{130}%
\begin{EESegment}{EE-TGT-P130-001}
\noindent\textbf{The axial pulse.} An axial pulse provides the adjustable contribution to
the total \NSi{NS-F-P130-001}{S}-moment. Two smaller corrections at the end of the
pulse set \NSi{NS-F-P130-002}{M=J=0}. Let \NSi{NS-F-P130-003}{X_{\mathrm p}} be the
radius at the end of the intermediate power-law interval, and write
\NSi{NS-F-P130-004}{E(X_{\mathrm p},\eta)=e_bf}. For
\NSi{NS-F-P130-005}{0\le y\le13/\lambda}, keep \NSi{NS-F-P130-006}{l=-\lambda}
and set \NSi{NS-F-P130-007}{U=ER_b}. With
\NSi{NS-F-P130-008}{\zeta_b=\lambda y}, define
\NSFormula{NS-F-P130-009}{display}{none}
\[
\begin{aligned}
\phi_b(\zeta_b)&=\int_0^{\zeta_b}\sigma(v/.02)\,dv
  &&(\zeta_b\ge0),
&R_0(\zeta_b)&=\phi_b(\zeta_b)[1-\sigma(\zeta_b-10)],\\
R_b&=\operatorname{Amp}(\eta)R_0(\zeta_b)+c_1\beta_1(y)+c_2\beta_2(y).
\end{aligned}
\]
Extend \NSi{NS-F-P130-010}{\phi_b} by zero when
\NSi{NS-F-P130-011}{\zeta_b<0}. The bumps \NSi{NS-F-P130-012}{\beta_1,\beta_2} are
identical translates of width \NSi{NS-F-P130-013}{.3}, centered at
\NSi{NS-F-P130-014}{13/\lambda-3} and \NSi{NS-F-P130-015}{13/\lambda-1}. Their
coefficients \NSi{NS-F-P130-016}{c_1,c_2} set \NSi{NS-F-P130-017}{M=J=0} at the
end of the pulse. These coefficients are affine functions of
\NSi{NS-F-P130-018}{\operatorname{Amp}}. The first equality in \eqref{eq:A.8} will
determine that amplitude.
\end{EESegment}

\begin{EESegment}{EE-TGT-P130-002}
\noindent\textbf{Profile interpolation and the angular moment correction.} Next, remove
the parameter dependence from the azimuthal profile and adjust its angular moment without
changing the pressure integral. Set \NSi{NS-F-P130-019}{U=0} from this point on. Let
\NSi{NS-F-P130-020}{X_{\mathrm{end}}} be the radius at the end of the pulse, and write
\NSi{NS-F-P130-021}{E(X_{\mathrm{end}},\eta)=e_{\mathrm{end}}f}. Set
\NSFormula{NS-F-P130-022}{display}{A.10}
\begin{equation}
\log E=\log e_{\mathrm{end}}-(1/2+\lambda)y-\vartheta_f(y)J_0
 -(1-\vartheta_f(y))\log2,
\qquad
\vartheta_f(y)=1-\sigma(y/T_f).
\tag{A.10}\label{eq:A.10}
\end{equation}
Choose a fixed, sufficiently large \NSi{NS-F-P130-023}{T_f}. Then
\NSi{NS-F-P130-024}{-\lambda-.1\le l\le-\lambda}. The factor that depends on
\NSi{NS-F-P130-025}{\eta} decreases to the constant
\NSi{NS-F-P130-026}{1/2}, so the amplitude cannot grow. Next, keep the resulting
\NSi{NS-F-P130-027}{\eta}-independent exponential for
\NSi{NS-F-P130-028}{30\log(1/\lambda)}. Place two relative
\NSi{NS-F-P130-029}{E} bumps of width \NSi{NS-F-P130-030}{.3}, centered three units
and one unit before the end of this interval. They impose
\NSFormula{NS-F-P130-031}{display}{A.11}
\begin{equation}
I=\frac{XH}{1-\lambda}\quad\text{at the end},
\qquad
\int_{\text{two bumps}}(E^2-E_{\mathrm{unedited}}^2)\,dy=0.
\tag{A.11}\label{eq:A.11}
\end{equation}
The value of \NSi{NS-F-P130-032}{E} at both endpoints stays unchanged. This part of
the construction uses only \NSi{NS-F-P130-033}{E} and does not depend on
\NSi{NS-F-P130-034}{\operatorname{Amp}}.
\end{EESegment}

\begin{EESegment}{EE-TGT-P130-003}
\noindent\textbf{Transition to the exterior power law.} Let
\NSi{NS-F-P130-035}{X_{\mathrm{rel}}} be the radius where the exterior transition starts,
and let \NSi{NS-F-P130-036}{e_{\mathrm{rel}}} be the unchanged, uniform value of
\NSi{NS-F-P130-037}{E} there. Over a unit interval, vary
\NSi{NS-F-P130-038}{l} smoothly from \NSi{NS-F-P130-039}{-\lambda} to
\NSi{NS-F-P130-040}{-1}. Keep \NSi{NS-F-P130-041}{l=-1} for
\NSi{NS-F-P130-042}{4\log(1/h)}, then vary it from \NSi{NS-F-P130-043}{-1} to
\NSi{NS-F-P130-044}{-h} over another unit interval. Each transition uses
\NSi{NS-F-P130-045}{\sigma}. On the last radial interval
\NSi{NS-F-P130-046}{0\le y\le3}, set
\NSFormula{NS-F-P130-047}{display}{A.12}
\begin{equation}
\psi_o(y)=1-\sigma((y-1)/2),
\qquad
f_o(y)=1-\rho_o\psi_o(y),
\qquad
\rho_o=c_oh,
\qquad
l=-h+\frac{f_o'}{f_o},
\tag{A.12}\label{eq:A.12}
\end{equation}
where \NSi{NS-F-P130-048}{c_o>0} is fixed small enough that
\NSi{NS-F-P130-049}{0\le f_o'/f_o<h/4}. Extend \NSi{NS-F-P130-050}{f_o} as a
constant to the left of this interval, and extend it by \NSi{NS-F-P130-051}{1} to the
right. Define
\NSFormula{NS-F-P130-052}{display}{A.13}
\begin{equation}
Q_p=\int_0^3\exp\left(\int_0^v(1+l(s))\,ds\right)\frac{f_o'(v)}{f_o(v)}\,dv.
\tag{A.13}\label{eq:A.13}
\end{equation}
After the transition to \NSi{NS-F-P130-053}{-h}, keep
\NSi{NS-F-P130-054}{l=-h} until the solution of
\NSi{NS-F-P130-055}{Q'+(1+l)Q=-l-h}, which starts at
\NSi{NS-F-P130-056}{Q=(\lambda-h)/(1-\lambda)} when the exterior transition begins,
has decreased to \NSi{NS-F-P130-057}{Q_p}. Run the final transition over three units of
logarithmic radial length, then keep \NSi{NS-F-P130-058}{l=-h} at every larger radius.
On and after that interval, \NSi{NS-F-P130-059}{E=E_{\mathrm{pow}}f_o}.
\end{EESegment}

\begin{proposition}\label{prop:A.4}
\begin{EESegment}{EE-TGT-P130-004}
There are finite choices, made in the order \eqref{eq:A.6}, for which the profile from
Section~\ref{sec:A.2} defines smooth \NSi{NS-F-P130-060}{U,E} for
\NSi{NS-F-P130-061}{X>0} and \NSi{NS-F-P130-062}{\eta\in[-1,1]}, with
\NSi{NS-F-P130-063}{E>0}, and satisfies \eqref{eq:A.8}. Every stage boundary divided
by \NSi{NS-F-P130-064}{X_R} is independent of both
\NSi{NS-F-P130-065}{X_R} and \NSi{NS-F-P130-066}{\eta}. After the compact axial
pulse, \NSi{NS-F-P130-067}{U=M=J=0}. After the final interval,
\NSi{NS-F-P130-068}{E=E_{\mathrm{pow}}}. There are four disjoint, unchanged patches,
listed after \eqref{eq:A.9}, where \NSi{NS-F-P130-069}{U=0} and
\NSi{NS-F-P130-070}{E} is a positive multiple of
\NSi{NS-F-P130-071}{fX^{-1/2-\lambda}}. For every fixed
\end{EESegment}
\NSPage{131}%
\begin{EESegment}{EE-TGT-P131-001}
\NSi{NS-F-P131-001}{x_-\in(0,1]}, a sufficiently large
\NSi{NS-F-P131-002}{X_R}, which may depend on \NSi{NS-F-P131-003}{x_-}, makes the
reference profile meet the strict relaxed cone condition from
\NSi{NS-F-P131-004}{x=x_-} through the first half-unit of its final tail. It meets the
admissible stress-cone condition throughout the
intermediate power-law interval, the axial pulse, the profile interpolation, and the exterior
transition. It also meets that condition as soon as \NSi{NS-F-P131-005}{l<0} during the
transition into the constant-slope interval. These statements cover a finite logarithmic
interval ending at tail coordinate \NSi{NS-F-P131-006}{1/2}. The heat construction below
handles the rest of the endpoint collar.
\end{EESegment}
\end{proposition}

\begin{EESegment}{EE-TGT-P131-002}
\subsection{Closing the moment integrals.}\label{sec:A.3}
Section~\ref{sec:A.2} fixed the radial pieces but left the amplitudes in the axial pulse
and the moment-correction bumps open. This subsection chooses those amplitudes so that the
moment identities in Proposition~\ref{prop:A.4} hold. The same estimates control the size
of the corrections and their parameter derivatives, which will be needed to check the
stress-cone inequalities in Section~\ref{sec:A.5}.
\end{EESegment}

\begin{EESegment}{EE-TGT-P131-003}
Write \NSi{NS-F-P131-007}{C_{\mathrm{pre}}} for constants that may depend on the fixed
choices \NSi{NS-F-P131-008}{M_d,T_d,P_*} and on the fixed steps. They are uniform for
small \NSi{NS-F-P131-009}{\lambda}, followed by sufficiently small
\NSi{NS-F-P131-010}{h}. Taking any prescribed finite number of parameter derivatives
may change these constants. None of them depends on the later radius
\NSi{NS-F-P131-011}{X_R} or on the normalization
\NSi{NS-F-P131-012}{\mathsf C}.
\end{EESegment}

\begin{EESegment}{EE-TGT-P131-004}
Directly integrating the axial transition and the intermediate power-law interval gives, at
the start of the pulse,
\NSFormula{NS-F-P131-013}{display}{A.14}
\begin{equation}
e_b\le C_{\mathrm{pre}}\lambda^{30},
\qquad
\left\lVert\frac{\mathcal A_X(U)}{E}\right\rVert_{C_\eta^1}
 \le C_{\mathrm{pre}}\lambda^{29},
\qquad
\left\lVert\frac{S(X_{\mathrm p})}{X_{\mathrm p}e_b^2f^2}\right\rVert_{C_\eta^1}
 \le C_{\mathrm{pre}}(1+\log(1/\lambda)).
\tag{A.14}\label{eq:A.14}
\end{equation}
Here is why these bounds hold. The moment \NSi{NS-F-P131-014}{M} stays constant after
the axial transition, while \NSi{NS-F-P131-015}{XE} grows at rate
\NSi{NS-F-P131-016}{1/2-\lambda} on the intermediate power-law interval. The
denominator \NSi{NS-F-P131-017}{XHE} for the
\NSi{NS-F-P131-018}{J} radial moment grows at rate
\NSi{NS-F-P131-019}{1/2-2\lambda}, which
gives the same kind of small bound for its moment discrepancy. Also,
\NSi{NS-F-P131-020}{XE^2} is constant during the axial transition and changes by the
factor \NSi{NS-F-P131-021}{e^{-2\lambda y}} on the intermediate power-law interval.
That factor stays bounded above and below for \NSi{NS-F-P131-022}{0\le y\le T_w}.
Integrating over the logarithmic length proves the last bound, including the fixed reference
profile on the inner interval. The function \NSi{NS-F-P131-023}{f} and its reciprocal
have bounded fixed derivatives on \NSi{NS-F-P131-024}{[-1,1]}, so differentiating these
integrals gives the stated parameter bounds.
\end{EESegment}

\begin{EESegment}{EE-TGT-P131-005}
For the two pulse corrections, divide by the matching normalizing factors
\NSi{NS-F-P131-025}{X_{\mathrm p}e_bf} and
\NSi{NS-F-P131-026}{X_{\mathrm p}H(X_{\mathrm p})e_bf} from
\NSi{NS-F-P131-027}{M,J}. In the \NSi{NS-F-P131-028}{dy} coordinate, the weights are
\NSi{NS-F-P131-029}{e^{s_1y},e^{s_2y}}, where
\NSi{NS-F-P131-030}{s_1=1/2-\lambda} and
\NSi{NS-F-P131-031}{s_2=1/2-2\lambda}. Normalize each weight at the first bump center.
The main pulse ends at \NSi{NS-F-P131-032}{11/\lambda}, which is at least
\NSi{NS-F-P131-033}{2/\lambda-3} before that center, and both slopes are larger than
\NSi{NS-F-P131-034}{.4}. The normalized moment discrepancy from the main pulse is
therefore at most
\NSi{NS-F-P131-035}{C(1+|\operatorname{Amp}|)e^{-c/\lambda}}. The contribution from
before the pulse picks up the factor
\NSi{NS-F-P131-036}{e^{-s_i(13/\lambda-3)}} in addition to the bounds just proved.
Formula~\eqref{eq:A.1} bounds the inverse matrix by
\NSi{NS-F-P131-037}{O(\lambda^{-1})}. The exponential absorbs that loss, giving
\NSFormula{NS-F-P131-038}{display}{A.15}
\begin{equation}
\lVert c_i\rVert_{C_\eta^1}
 \le C_{\mathrm{pre}}e^{-c/\lambda}(1+|\operatorname{Amp}|),
\qquad
M=J=0\quad\text{after the pulse}.
\tag{A.15}\label{eq:A.15}
\end{equation}
Both here and in the amplitude calculation, this estimate says that the affine coefficients
of \NSi{NS-F-P131-039}{1} and \NSi{NS-F-P131-040}{\operatorname{Amp}} have the
stated parameter bounds. Every fixed \NSi{NS-F-P131-041}{y} derivative has the same
bound. Every fixed \NSi{NS-F-P131-042}{\zeta_b} derivative brings in only a power of
\NSi{NS-F-P131-043}{\lambda^{-1}}, so it is still exponentially small.
\end{EESegment}

\begin{EESegment}{EE-TGT-P131-006}
The axial corrections now fix \NSi{NS-F-P131-044}{M} and
\NSi{NS-F-P131-045}{J} for each trial pulse amplitude. Next, set the angular moment
needed at the start of the exterior transition. Put
\NSi{NS-F-P131-046}{r_I=I/(XH)}. It satisfies
\NSFormula{NS-F-P131-047}{display}{none}
\[
r_I'+(1+l)r_I=1.
\]
Before the exterior transition, \NSi{NS-F-P131-048}{1+l} stays bounded below, and the
explicit shapes give a \NSi{NS-F-P131-049}{C_\eta^1} bound
\NSi{NS-F-P131-050}{C_{\mathrm{pre}}} at the end of the profile interpolation. On the
interval with uniform profile, \NSi{NS-F-P131-051}{r_I-}
\end{EESegment}
\NSPage{132}%
\begin{EESegment}{EE-TGT-P132-001}
\NSi{NS-F-P132-001}{(1-\lambda)^{-1}} decays at rate
\NSi{NS-F-P132-002}{1-\lambda}. The moment discrepancy at the first correction center
is therefore \NSi{NS-F-P132-003}{O_{C_\eta^1}(C_{\mathrm{pre}}\lambda^{28})}. Make a
relative change \NSi{NS-F-P132-004}{E\mapsto E(1+c_1\beta_1+c_2\beta_2)}. The
\NSi{NS-F-P132-005}{I} row is linear with slope
\NSi{NS-F-P132-006}{1-\lambda}, and the linearized pressure row has slope
\NSi{NS-F-P132-007}{-1-2\lambda} in \NSi{NS-F-P132-008}{y}. Divide these rows by
\NSi{NS-F-P132-009}{XH} and \NSi{NS-F-P132-010}{E^2}, respectively, at the first
center. Equation~\eqref{eq:A.1} gives an inverse that stays bounded. The quadratic
pressure increment is the only nonlinear term. Lemma~\ref{lem:A.2} gives coefficients
\NSi{NS-F-P132-011}{O_{C_\eta^1}(C_{\mathrm{pre}}\lambda^{28})} and the exact
conditions in \eqref{eq:A.11}. The \NSi{NS-F-P132-012}{y} derivatives of these
corrections are much smaller than \NSi{NS-F-P132-013}{\lambda}, so
\NSi{NS-F-P132-014}{E>0} and \NSi{NS-F-P132-015}{l\le-\lambda/2} both remain true.
\end{EESegment}

\begin{EESegment}{EE-TGT-P132-002}
At the start of the exterior transition,
\NSi{NS-F-P132-016}{U=M=J=0} and
\NSi{NS-F-P132-017}{I=XH/(1-\lambda)} does not depend on
\NSi{NS-F-P132-018}{\eta}. The integral formula \eqref{eq:4.16} for
\NSi{NS-F-P132-019}{Q_s} gives
\NSi{NS-F-P132-020}{Q_s=(\lambda-h)/(1-\lambda)} exactly. The source
\NSi{NS-F-P132-021}{-l-h} is nonnegative from there until the final interval. During the
first transition, that solution keeps the lower bound
\NSi{NS-F-P132-022}{c\lambda}. On
the constant-slope interval with
\NSi{NS-F-P132-023}{l=-1}, \NSi{NS-F-P132-024}{Q_s'=1-h}, so the value at the end is
comparable to \NSi{NS-F-P132-025}{\log(1/h)}. The next transition changes that size by
only bounded factors. By contrast, \NSi{NS-F-P132-026}{Q_p\asymp\rho_o\asymp h}, as
\eqref{eq:A.13} shows on a fixed interval. The constant-slope interval in between, where
\NSi{NS-F-P132-027}{l=-h}, therefore has a positive, finite logarithmic radial length.
On the final interval, the exact solution is
\NSFormula{NS-F-P132-028}{display}{A.16}
\begin{equation}
Q_s(y)=\int_y^3\exp\left(\int_y^v(1+l(s))\,ds\right)\frac{f_o'(v)}{f_o(v)}\,dv.
\tag{A.16}\label{eq:A.16}
\end{equation}
It reaches zero at the endpoint and stays zero afterward. Throughout the exterior
transition, \NSi{NS-F-P132-029}{I} does not depend on
\NSi{NS-F-P132-030}{\eta}, and after that
\NSi{NS-F-P132-031}{Q_s=-1+(1-h)I/(XH)}. It follows that
\NSFormula{NS-F-P132-032}{display}{none}
\[
I=\frac{XH}{1-h}\quad\text{beyond the tail},
\qquad
\int_0^\infty(H-H_{\mathrm{pow}})\,dX=0.
\]
\end{EESegment}

\begin{EESegment}{EE-TGT-P132-003}
The subtracted moment is integrable at zero because
\NSi{NS-F-P132-033}{H_{\mathrm{pow}}\propto X^{-h}} and
\NSi{NS-F-P132-034}{h<1}. Far enough out, its integrand is identically zero. The
angular correction, its endpoint, and all the lengths in the exterior transition depend only
on \NSi{NS-F-P132-035}{E}, so they do not depend on the axial amplitude.
\end{EESegment}

\begin{EESegment}{EE-TGT-P132-004}
The angular moment has now been fixed independently of the pulse amplitude. The one scalar
condition left is \NSi{NS-F-P132-036}{S(\infty)=0}. Before solving it, bound the part
that comes from the transition to the exterior profile. Since
\NSi{NS-F-P132-037}{(XE^2)'=2lXE^2}, the interval with constant slope
\NSi{NS-F-P132-038}{l=-1} gives the exact factors
\NSFormula{NS-F-P132-039}{display}{A.17}
\begin{equation}
\frac{(XE^2)_{\mathrm{hold\ end}}}{(XE^2)_{\mathrm{hold\ start}}}=h^8,
\qquad
\frac{E_{\mathrm{hold\ end}}}{E_{\mathrm{hold\ start}}}=h^6.
\tag{A.17}\label{eq:A.17}
\end{equation}
Here ``hold start'' and ``hold end'' are the two endpoints of the constant-slope interval
where \NSi{NS-F-P132-040}{l=-1}. The label ``release'' in the next integral means the
whole transition to the exterior power law. The first transition and the constant-slope
interval contribute at most
\NSi{NS-F-P132-041}{CX_{\mathrm{rel}}e_{\mathrm{rel}}^2} to
\NSi{NS-F-P132-042}{\int XE^2\,dy}. From the end of the constant-slope interval onward,
\NSi{NS-F-P132-043}{l\le-3h/4}, so the rest of the integral is at most
\NSi{NS-F-P132-044}{CX_{\mathrm{rel}}e_{\mathrm{rel}}^2h^8/h}. Therefore
\NSFormula{NS-F-P132-045}{display}{A.18}
\begin{equation}
\int_{\mathrm{release}}XE^2\,dy\le CX_{\mathrm{rel}}e_{\mathrm{rel}}^2
\tag{A.18}\label{eq:A.18}
\end{equation}
with a constant that is uniform in \NSi{NS-F-P132-046}{h}. The
\NSi{NS-F-P132-047}{h^8} estimate controls this integral. The separate
\NSi{NS-F-P132-048}{h^6} estimate will later control the ratio
\NSi{NS-F-P132-049}{w=N_s/(EQ_s)} near the outer endpoint. The profile interpolation
and the constant-slope interval after it contribute at most
\NSi{NS-F-P132-050}{C(1+\log(1/\lambda))X_{\mathrm{end}}e_{\mathrm{end}}^2}. The first
\NSi{NS-F-P132-051}{\eta} derivatives satisfy the same bounds because the exterior
transition is uniform and the earlier logarithmic shapes have bounded derivatives.
\end{EESegment}
\par
